\section*{Capítulo \ref{ch:g5:natural-hazards} --- Los peligros naturales}

\begin{solution}{exo:g5:natural-hazards:1}
(a) Solo un \omterm{def:g5:natural-hazards:hazard}{peligro natural}: no hay daños para nadie. (b) Una \omterm{def:g5:natural-hazards:disaster}{catástrofe
natural}: se dañan viviendas. (c) Solo un \omterm{def:g5:natural-hazards:hazard}{peligro natural}. (d) Una
\omterm{def:g5:natural-hazards:disaster}{catástrofe natural}.
\end{solution}

\begin{solution}{exo:g5:natural-hazards:2}
Del aire: una ola de calor, una \omterm{def:g5:natural-hazards:storm}{tormenta}, una \omterm{def:g5:natural-hazards:drought}{sequía}. Del agua: una
\omterm{def:g5:water-shapes-the-land:flood}{inundación}. Del suelo: un \omterm{def:g5:natural-hazards:landslide}{deslizamiento de tierra}, una \omterm{def:g4:volcanoes-earthquakes-first-look:eruption}{erupción} volcánica,
un \omterm{def:g4:volcanoes-earthquakes-first-look:earthquake}{terremoto}.
\end{solution}

\begin{solution}{exo:g5:natural-hazards:3}
Una \omterm{def:g5:natural-hazards:drought}{sequía} es un periodo largo, de meses o de años, con mucha menos \omterm{def:g4:clouds-rain-snow:precipitation}{lluvia}
de la habitual. Puede dejar a los agricultores sin cosecha, porque los
cultivos se marchitan y mueren, y puede dejar a las ciudades sin agua
potable suficiente, porque los \omterm{def:g1:landscapes-around-us:water}{ríos} y los \omterm{def:g1:landscapes-around-us:water}{lagos} menguan.
\end{solution}

\begin{solution}{exo:g5:natural-hazards:4}
Por ejemplo: meter dentro o atar todo lo que el viento pueda llevarse
(sillas de jardín, juguetes, cubos de basura); comprobar las linternas y
tener preparadas pilas de repuesto y agua; quedarse dentro, lejos de las
ventanas, durante la \omterm{def:g5:natural-hazards:storm}{tormenta}; seguir escuchando las noticias en la radio.
\end{solution}

\begin{solution}{exo:g5:natural-hazards:5}
(a) Un huracán. (b) Un tifón. (c) Un ciclón tropical.
\end{solution}

\begin{solution}{exo:g5:natural-hazards:6}
Verde, amarillo, naranja, rojo. En el nivel rojo hay un gran peligro: hay
que actuar enseguida y seguir las instrucciones de los responsables.
\end{solution}

\begin{solution}{exo:g5:natural-hazards:7}
Una parte de la ladera de la derecha se desprendió y se deslizó hacia abajo,
dejando una cicatriz desnuda. Su \omterm{def:g3:stones-pebbles-sand:sand}{barro} se extendió por todo el fondo del
\omterm{def:g1:landscapes-around-us:landform}{valle} y cubrió la carretera. También cortó el \omterm{def:g1:landscapes-around-us:water}{río} como una presa: el agua
del \omterm{def:g1:landscapes-around-us:water}{río}, que ya no puede seguir su camino, se ha acumulado detrás y ha
inundado el bosque de \omterm{def:g1:landscapes-around-us:water}{río} arriba.
\end{solution}

\begin{solution}{exo:g5:natural-hazards:8}
Lena tiene razón. No podemos detener la \omterm{def:g4:clouds-rain-snow:precipitation}{lluvia} ni el \omterm{def:g1:landscapes-around-us:water}{río}, pero sí podemos
poner a la gente a salvo. Por ejemplo: construir las casas en terreno alto
en lugar de en la \omterm{def:g5:water-shapes-the-land:flood}{llanura de inundación}; avisar a tiempo a la gente con
avisos de \omterm{def:g5:water-shapes-the-land:flood}{inundación}, para que suba a terreno más alto; enseñar a todos a
no atravesar nunca el agua de una \omterm{def:g5:water-shapes-the-land:flood}{inundación} a pie ni en coche.
\end{solution}

\begin{solution}{exo:g5:natural-hazards:9}
(a) $4 \times 60 = 240$ milímetros. (b) La cuarta parte de 60 es 15; en
cuatro meses, $4 \times 15 = 60$ milímetros. (c) Faltaron
$240 - 60 = 180$ milímetros: tres cuartas partes de la \omterm{def:g4:clouds-rain-snow:precipitation}{lluvia} habitual.
\end{solution}

\begin{solution}{exo:g5:natural-hazards:10}
(a) $1\,200 \div 4 = 300$ viviendas. (b) Un tercio de 300 es 100; así que
siguen en la \omterm{def:g5:water-shapes-the-land:flood}{llanura de inundación} $300 - 100 = 200$ viviendas.
\end{solution}

\begin{solution}{exo:g5:natural-hazards:11}
Razones posibles: el primer pueblo recibió el aviso a tiempo y la gente se
preparó (contraventanas cerradas, objetos metidos dentro), y el segundo no;
las casas del primero están bien construidas, y las del segundo son
frágiles; el segundo pueblo está en una \omterm{def:g1:landscapes-around-us:water}{costa} baja o en una \omterm{def:g5:water-shapes-the-land:flood}{llanura de
inundación}, donde la \omterm{def:g5:natural-hazards:storm}{tormenta} trajo además una \omterm{def:g5:water-shapes-the-land:flood}{inundación}, mientras que el
primero está en terreno alto; en el segundo pueblo, los árboles cercanos a
las casas cayeron sobre ellas. El mismo peligro no causa el mismo daño en
todas partes: depende de dónde y cómo vive la gente, y de lo bien preparada
que esté.
\end{solution}
